\section{Rescaling, contraction and matching}\label{sec:rescaling}

Let $f_t$ be the local actual-moduli comparison of
Theorem~\ref{thm:harmonic-comparison}. Write $g_t^{L^2}$ for the
actual harmonic metric and $g_{\zeta(t)}$ for the finite quotient metric.
We use the estimate
\begin{equation}\label{eq:actual-relative}
 (1-C(|t|+r^\alpha))g_{\zeta(t)}
 \le f_t^*g_t^{L^2}
 \le(1+C(|t|+r^\alpha))g_{\zeta(t)},\qquad r=\|x\|,
\end{equation}
where $\alpha=1$ in general and $\alpha=2$ under that theorem's
verified parity hypothesis. Actual distances use the corresponding
finite-stratum curve class in the full bundle-moduli space.
Global completeness of that space is not assumed.

\subsection{Localization and distances to wall strata}
Shrink the common chart so that it contains $r<2R$ and has coarse
squared-metric bounds $1/2$ and $2$. Each closed radial sublevel
is compact in the quotient topology, so its image is compact and
closed in the Hausdorff actual moduli space.
A path starting at radius $\rho<R$ and escaping the image of $r<R$
must first reach radius $R$. Its initial segment has length at
least $(R-\rho)/\sqrt2$, by the $1$-Lipschitz radius.
While a path of length $L$ stays in the chart, its lifted radius
is at most $\rho+\sqrt2L$. These observations exclude short paths
that leave the local model.

\begin{proposition}\label{prop:stratum-distance}
At $t=0$, let $\Sigma\subset M_0$ be a closed conical union of
representation-type strata containing the vertex. Suppose the
specified actual target has intersection with the chart equal to
$f_0(\Sigma)$. Then, as $r(x)\to0$,
\[
 d^{L^2}_0(f_0(x),f_0(\Sigma))
  =(1+O(r(x)^\alpha))d_0(x,\Sigma).
\]
For the vertex this is
$d^{L^2}_0(f_0(x),[E_0])=r(x)(1+O(r(x)^\alpha))$.
\end{proposition}
\begin{proof}
The finite distance to $\Sigma$ is at most $r(x)$, using the vertex.
An admissible near-minimizer of length at most $d_0(x,\Sigma)+\delta$
stays within radius $2r(x)+\delta$. Its actual image is a competitor.
Any shorter actual competitor has length $O(r(x))$, cannot escape
the larger chart, and stays within radius $C_0r(x)$ by the preceding
barrier. Its endpoint cannot lie in a portion of the target outside
the chart without paying that exit cost.
Apply \eqref{eq:actual-relative} on these paths and let $\delta\downarrow0$.
The estimate is multiplicative even when the target distance is much
smaller than $r(x)$. The vertex formula uses \eqref{eq:cone-radius}.
\end{proof}
The target identification is a geometric hypothesis: a tensor estimate
alone does not identify a named singular stratum.

\subsection{Arbitrary directions and complete actual limits}
Assume the finite complex zero fiber contains a module-simple point,
as in Section~\ref{sec:intrinsic}.
\begin{theorem}\label{thm:actual-limits}
Let $t_n\to0$ and $\delta_n\downarrow0$, with
$\eta_n=\zeta(t_n)/\delta_n^2\to\eta_*$.
Let balanced base representatives $p_n\in A_{\eta_n}$ converge to $p_*$.
With base point $f_{t_n}([\delta_np_n])$ and distances divided by
$\delta_n$, the actual moduli have complete pointed local limit
$(M_{\eta_*},d_{\eta_*},[p_*])$.
On each bounded rescaled ball, the intrinsic comparison with
$M_{\eta_n}$ is relative with error
$O_L(|t_n|+\delta_n^\alpha)$.
\end{theorem}
\begin{proof}
Homogeneity identifies $M_{\delta_n^2\eta_n}$ with
$\delta_nM_{\eta_n}$, multiplying distances by $\delta_n$.
The fixed chart has an exit cost bounded below independently of $n$
from starting radius $O(\delta_n)$. Any actual path of length
$O(\delta_n)$ stays inside it, and its finite radius is $O(\delta_n)$
by the coarse comparison.
Conversely, admissible finite near-minimizers between points in a
bounded model ball have uniformly bounded radial coordinate by
the $1$-Lipschitz radius and their length bound. They scale into
the actual chart. On both types of paths \eqref{eq:actual-relative}
has the claimed error. Taking infima gives both intrinsic inequalities,
including pairs at distances arbitrarily smaller than $\delta_n$.

Full inverse coverage puts every actual radius-$\delta_nL$ ball
inside the image of a model radius-$L+o(1)$ ball, and gives the
reverse inclusion with radius $L-o(1)$.
The model balls are compact; local distance comparison makes these
actual balls compact as well. Theorem~\ref{thm:varying-parameters}
now proves convergence to the entire complete target. No global
properness of the actual bundle-moduli space is used.
\end{proof}

Along a fixed ray $\zeta(t_\kappa)=\kappa\eta$ this gives the scale
$\sqrt\kappa$ and error $O_L(|t_\kappa|+\kappa^{\alpha/2})$.
For any specified compact subset $C\subset M_\eta$ of positive diameter,
\begin{equation}\label{eq:fixed-core-scale}
 \operatorname{diam}_{L^2}
 f_{t_\kappa}(\sqrt\kappa C)
 =\sqrt\kappa\bigl(1+O(|t_\kappa|+\kappa^{\alpha/2})\bigr)
      \operatorname{diam}_{d_\eta}C.
\end{equation}
Identifying $C$ with a contracted set is a separate geometric step.

\subsection{Normal representations and moving centers}
\begin{lemma}\label{lem:normal-simple}
If the complex zero fiber contains a module-simple representation,
the full normal Ext-quiver at
every point of every $M_\eta$ also has a module-simple representation
in its complex zero fiber.
\end{lemma}
\begin{proof}
The common module-simple locus gives a dense open subset of the
complex quotient, hence of $M_\eta$ by
Corollary~\ref{cor:fixed-face}. Choose balanced representatives
$x_j\to x_*$ from it. The common-parameter inverse in
Theorem~\ref{thm:metric-slice}, at fixed parameter $\eta$, gives
$w_j\to0$ in the full normal $W$ with both normal moments zero.
The full equivariant slice and its complex corrections identify
stabilizers up to conjugacy. Thus $H_{w_j}$ is exactly the single
common scalar circle, since this is the stabilizer of $x_j$.
Proposition~\ref{prop:real-polystability} complexifies it.
At zero level, that proposition decomposes the representation
into module-simple summands: zero-stability means that there are
no proper subrepresentations. More than one summand or a repeated
summand would give a nonscalar endomorphism. Hence $w_j$ is
module-simple.
\end{proof}

Apply all quiver statements to the full $W$, not to $W'$ in
\eqref{eq:normal-product}. Properness and parameter convergence
for $W'///_\nu H$ follow from the exact product and restriction
to points with $U$-coordinate zero.

\begin{theorem}\label{thm:moving-centers}
Let $\eta_n\to\eta_*$ and $p_n\in M_{\eta_n}$ approach $[x_*]$.
In the common slice of Theorem~\ref{thm:metric-slice}, write their
balanced normal coordinates as $(u_n,z_n)\to0$ in $U\oplus W'$.
Set
\[
 \zeta_{H,n}=\eta_n|_{\mathfrak h}
 =(\eta_n-\eta_*)|_{\mathfrak h}.
\]
For any $\sigma_n\downarrow0$ such that
\[
 \zeta_{H,n}/\sigma_n^2\to\nu,\qquad z_n/\sigma_n\to z
\]
after unitary adjustment, the complete pointed limit with distances
$d_{\eta_n}/\sigma_n$ and base $p_n$ is
\[
 U\times(W'///_\nu H),\qquad\text{with base }(0,[z]).
\]
The same conclusion holds inside the actual family of
Theorem~\ref{thm:actual-limits}, at physical scale $\delta_n\sigma_n$.
\end{theorem}
\begin{proof}
The exact model coordinate change is
\[
 (u,z')\longmapsto(u_n+\sigma_nu,\sigma_nz').
\]
Translation in $U$ does not affect the moments. It identifies the
rescaled normal model with $U\times(W'///_{\nu_n}H)$, where
$\nu_n=\zeta_{H,n}/\sigma_n^2$.
On a fixed bounded radial region its relative squared-metric error
in the original finite space is
\begin{equation}\label{eq:moving-error}
 O_{x_*}(|\eta_n-\eta_*|+|u_n|+\sigma_nR)=o(1).
\end{equation}
In particular this error is not divided by $\sigma_n$.

Choose a fixed radius $\rho$ whose twice-sized sublevel lies in
the slice. Compact sublevels and the coarse metric comparison give
exit cost at least $(\rho-\|(u_n,z_n)\|)/\sqrt2$.
After division by $\sigma_n$ this tends to infinity.
Every competing path of length $O(\sigma_n)$ remains in the chart,
and the radius of its translated coordinate $(u-u_n,z')$ is bounded
by $\|z_n\|+\sqrt2$ times its length, hence is $O(\sigma_n)$.
Conversely, admissible normal-model near-minimizers on bounded
rescaled balls stay in bounded translated radial regions.
Equation~\eqref{eq:moving-error} therefore compares both intrinsic
infima relatively; the common inverse supplies full ball coverage.
Lemma~\ref{lem:normal-simple} and
Theorem~\ref{thm:varying-parameters}, applied to full $W$, give the
whole complete product limit.

In the actual family the outer error is
$O(|t_n|+\delta_n^\alpha R^\alpha)$ on a bounded region of the
first rescaled model. The inner error is \eqref{eq:moving-error}.
Both are multiplicative, so they add to first order and are not
amplified by the smaller scale. The outer chart and the fixed inner
slice each have diverging exit cost after rescaling by
$\delta_n\sigma_n$. The same inverse-coverage and path argument
therefore proves the actual conclusion.
\end{proof}

Tangential motion $u_n$ may be arbitrarily large compared with
$\sigma_n$ while tending to zero. There is no comparability
assumption between independent stability parameters.
For example, when nonzero,
$\sigma_n=\max(\|z_n\|,\sqrt{\|\zeta_{H,n}\|})$ gives bounded
rescaled normal parameters and representatives, hence convergent
subsequences. Smaller scales are allowed whenever the stated
boundedness conditions hold.

\begin{corollary}\label{cor:iterated-limits}
Theorem~\ref{thm:moving-centers} can be iterated through any fixed
finite succession of convergent centers and positive scales tending
to zero. At each stage require convergence of the retained parameter
divided by the square of its scale and of the rescaled normal base
representatives. The physical scale is the product of all stage
scales, and the complete limit is the corresponding full normal
quotient with its trivial Euclidean factors.
\end{corollary}
\begin{proof}
Lemma~\ref{lem:normal-simple} preserves the existence of a module-simple point in the complex zero fiber
at every stage. Apply the preceding theorem at each fixed limiting
center, with that center's own constants and radius.
For finitely many stages the relative errors sum to $o(1)$;
successive first-exit bounds and inverse charts give full ball
coverage. No uniform radius over all possible centers or convergence
of an infinite tower is asserted.
\end{proof}

\subsection{The size of the affine contracted fiber}
Let $C_\eta$ be the fiber over the affine vertex under
$Q_F\to Q_0$. It is nonempty and compact by the surjective projective
face map.
\begin{corollary}\label{cor:uniform-core}
There is $C$ such that
$\operatorname{diam}_{d_\eta}C_\eta\le C\sqrt{\|\eta\|}$
for every central $\eta$.
The same upper bound holds for its actual image $f_t(C_{\zeta(t)})$
when $\zeta(t)\to0$ and $t\to0$.
\end{corollary}
\begin{proof}
Lemma~\ref{lem:invariant-properness} bounds representative radii
uniformly on $\|\eta\|=1$.
If the corresponding intrinsic diameters were unbounded, choose
endpoints and parameters, pass to convergent representatives, and
apply \eqref{eq:distance-continuity}, a contradiction.
Homogeneity gives the asserted scale, including the zero fiber,
which is the vertex. The actual assertion follows by localization
and \eqref{eq:actual-relative}.
\end{proof}
There is no uniform lower diameter estimate in this argument;
secondary contractions can occur on smaller scales.
The actual statement concerns the displayed image of the finite
core. A global contraction is not identified by this estimate alone.

\subsection{Matching to the surrounding zero-level metric}
At a zero-level point $x$ with scalar stabilizer, let $\lambda(x)>0$
be the least eigenvalue of $C_x^*C_x$ on the total-trace-zero
Hermitian vertex algebra, using its Hilbert--Schmidt norm.
\begin{proposition}\label{prop:outer-matching}
If $\|\zeta\|/\lambda(x)$ is sufficiently small, the same complex
orbit has a $\zeta$-balanced representative $e^Sx$ with
\[
 \|S\|_{\rm op}\le A\frac{\|\zeta\|}{\lambda(x)}.
\]
On the corresponding free quotient loci,
\[
 g_\zeta=\left(1+O\left(\frac{\|\zeta\|}{\lambda(x)}\right)\right)g_0
\]
as two-sided Hermitian-form inequalities.
The relative comparison with the actual metric has error
$O(|t|+\|x\|^\alpha+\|\zeta\|/\lambda(x))$.
\end{proposition}
\begin{proof}
Use the degree homotopy $\mu(e^Sx)-s\zeta$ on the operator-norm
ball $a=A\|\zeta\|/\lambda(x)$.
For a strict spectral cut $P$ put
$P_0=P-(\operatorname{tr}P/M)I$. Then
\[
 \|P_0\|^2=\frac{k(M-k)}M\ge\frac{M-1}M,\qquad
 q_P(x)\ge\lambda(x)\|P_0\|^2.
\]
The boundary estimate \eqref{eq:cut-lower} is at least
$2e^{-4a}a\lambda(x)\|P_0\|^2/(M-1)$, whereas the target pairing
has absolute value at most $\|\zeta\|\|P_0\|$.
Choose $A$ large depending only on $M$ and then $a$ small.
The homotopy has no boundary zeros and initial degree $+1$, giving
the correction. No stronger inverse-function condition involving
$\|\zeta\|\|x\|^2/\lambda(x)^2$ is used.

A zero-balanced point with scalar stabilizer is module-simple:
\eqref{eq:subrepresentation-trace} makes any invariant subspace
reducing and hence produces a nonscalar stabilizer if proper.
Thus these orbits belong to the common stable open and determine
the same affine quotient point.
On local holomorphic quotient tangents the derivative of the
complex correction is vertical. Removing it by the full complex
orbit infimum leaves only the action of $e^S$ on ambient norms,
whose squared distortion is $e^{\pm4\|S\|_{\rm op}}$.
This proves the finite metric estimate. Compose with
\eqref{eq:actual-relative} for the actual estimate.
\end{proof}
The matching statement applies when $\|\zeta\|/\lambda(x)\to0$.
The complete normal models above describe regimes where this
outer condition fails. None of these local stationary results
implies bubbling control or an HYM evolution theorem.

